\section*{4.4 Gadget classification (author source pp.22--23)}
Moreover, for conjunction DAGs we can prove an analog of Theorem 3.4:
\begin{originalresult}{Theorem 4.2}
Let $f\subseteq\bits\times V$ be a relation and $g:\bits^m\to\bits$ be a weakly resistant gadget. If there is a conjunction DAG for $f\comp g$ with $S$ leaves, then the certificate complexity of $f$ is at most $O(\log S)$.
\end{originalresult}
\begin{proof}
The proof of Theorem 3.4 can be applied in this setting as is because it ignores the internal structure of the decision tree and considers only the leaves.
\end{proof}
Note that Theorem 4.2 gives us a lifting from certificate complexity to leaf size. Unfortunately, in the case of the falsified clause problem for CNF, this theorem cannot be effectively used to prove lower bounds since leaf size, as well as certificate complexity, is usually small. However, this theorem still implies a dichotomy result similar to one in Theorem 3.9.
\begin{originalresult}{Theorem 4.3}
For every gadget $g:\bits^m\to\bits$, one of the following cases holds:
\begin{enumerate}
\item[(i)] There is an infinite family of functions $f_n$ such that $f_n\comp g$ can be computed by a conjunction DAG having $O(n)$ leaves, but $\CC(f_n)=\Omega(n)$.
\item[(ii)] For every relation $f$, if we have a conjunction DAG for $f\comp g$ with $S$ leaves, then the certificate complexity of $f$ is at most $O(\log S)$.
\end{enumerate}
\end{originalresult}
\begin{proof}
The two cases of the theorem correspond to the two cases of Lemma 3.8. If $g$ is a disjunction then we take $f_n:=\OR_n$. Clearly, $\CC(f_n)=n$ and $f_n\comp g$ can be computed with a conjunction DAG (actually a tree) having $O(n)$ leaves (if $g$ is constant, then the DAG has only one leaf). The case of a conjunction is similar.
Otherwise, gadget $g$ is weakly resistant and by Theorem 4.2 we have $\log S=\Omega(\CC(f))$.
\end{proof}
